\documentclass[11pt]{article}
\usepackage[margin=1in]{geometry}
\usepackage{booktabs}
\usepackage{longtable}
\usepackage{hyperref}
\usepackage{xcolor}
\usepackage{enumitem}
\usepackage{parskip}

\title{Replication Report: Shrestha et al.\ (2022)\\
\emph{Complete Genome Sequence and Comparative Genome Analysis of Variovorax sp.\ Strains PAMC28711, PAMC26660, and PAMC28562 and Trehalose Metabolic Pathways in Antarctica Isolates}}
\author{Ollie (OpenClaw AI) \\ BVBRC-47 \textperiodcentered\ TOPUP85 rank-27}
\date{2026-07-01}

\begin{document}
\maketitle

\begin{center}
\fbox{\textbf{Verdict: REPLICATED}}\\[2pt]
\small Coverage 8/10 \textperiodcentered\ Agreement 9/10 \textperiodcentered\ Independent free LLM-judge (Argo \texttt{gpt-5.2}): REPLICATED
\end{center}

\textbf{Paper:} Shrestha P, et al.\ \emph{International Journal of Microbiology} 2022, Art.\ 5067074.\\
\textbf{DOI:} \href{https://doi.org/10.1155/2022/5067074}{10.1155/2022/5067074} \textperiodcentered\ \textbf{PMC:} PMC10232917 \textperiodcentered\ \textbf{PMID:} 37275508 \textperiodcentered\ Open access (CC BY).

\paragraph{Dedup note.}
This is \emph{not} sibling directory \texttt{BVBRC-04-Variovorax-trehalose-Shrestha2022}. BVBRC-04 = Shrestha 2022 \emph{BMC Genomic Data} 23:4 (DOI 10.1186/s12863-021-01020-y), a single-strain (PAMC28711) trehalose-pathway prediction methods paper. BVBRC-47 (this report) = Shrestha 2022 \emph{Int J Microbiology} (DOI 10.1155/2022/5067074), a three-strain complete-genome + comparative-genomics paper. Different journal, DOI, PMC, and scope.

\section{Paper summary}

The paper reports the complete PacBio SMRT genomes of three Antarctic \emph{Variovorax} isolates
(PAMC28711, PAMC26660, PAMC28562, from Barton Peninsula, King George Island) and performs a
comparative-genomics survey against 16 other complete \emph{Variovorax} genomes (19 total in Table~1).
It covers: (i) genome statistics (size, GC\%, CDS, genes, tRNAs), (ii) whole-genome taxonomy via
TYGS + ANI (OAT/ANIb, ANIm) + digital DDH (GGDC), (iii) CAZyme content via the dbCAN2 meta server
(Table~3), and (iv) the trehalose biosynthesis/degradation gene inventory (OtsA/OtsB, TreS,
TreY/TreZ synthesis routes; trehalases). The headline biological finding is that
\textbf{PAMC28711 and PAMC28562 each carry all three trehalose biosynthetic pathways;
PAMC26660 carries only OtsA/OtsB}, framed as a cold/osmotic-stress adaptation. The three PAMC
strains have the lowest GC\% of the 19 compared genomes.

\section{Claims tested}

\begin{longtable}{@{}clllcc@{}}
\toprule
\# & Claim (short) & Type & Testable from public artifacts? & Tested \\
\midrule
C1 & Genome sizes 4.32/7.39/4.69 Mb; single circular chromosomes & Genome stat & Yes (RefSeq) & \checkmark \\
C2 & GC\% 66.00/66.00/63.70; lowest of 19 & Genome stat & Yes & \checkmark \\
C3 & tRNA counts 46/52/47 & Genome stat & Yes (RefSeq GFF) & \checkmark \\
C4 & Gene counts $\approx$4232/6919/4402 & Genome stat & Yes & \checkmark \\
\textbf{C5} & \textbf{PAMC28711 \& PAMC28562 have 3 trehalose pathways; PAMC26660 has 1 (OtsA/OtsB only)} & \textbf{Headline biology} & \textbf{Yes} & \checkmark \\
C6 & PAMC strains distinct at species level from V.\ paradoxus etc.\ (ANI<95\%, dDDH<70\%) & Taxonomy & Yes (fastANI proxy) & \checkmark \\
C7 & Per-strain CAZyme family counts (Table~3) & CAZyme & Partial (needs dbCAN2) & partial \\
C8 & AZCL wet-lab polysaccharide screen (Table~5) & Wet-lab & No & N/A \\
\bottomrule
\end{longtable}

\section{Method}

All work on free endpoints (Europe PMC full-text XML, NCBI Datasets, local
\texttt{fastANI}+\texttt{BLAST+}+\texttt{prokka} env on \texttt{uicgpu01}, Argo \texttt{gpt-5.2}
LLM-judge). No paid services.

\begin{enumerate}[leftmargin=1.4em]
  \item \textbf{Paper acquisition.} Europe PMC \texttt{PMC10232917/fullTextXML}
        (162~KB, SHA-256 \texttt{63620a15\ldots d3c0}); extracted abstract, Materials \& Methods,
        and all 5 tables. Hindawi PDF endpoint returns a Cloudflare HTML block, so XML is canonical.
  \item \textbf{Accession resolution.} Mapped paper nucleotide accessions (CP014517;
        CP060295/NZ\_CP060295; CP060296/NZ\_CP060296) to RefSeq assemblies via
        NCBI \texttt{esearch}/\texttt{esummary}:
        PAMC28711 $\to$ GCF\_001577265.1;
        PAMC26660 $\to$ GCF\_014302995.1;
        PAMC28562 $\to$ GCF\_014303735.1;
        \emph{V.\ paradoxus} NBRC 15149\textsuperscript{T} $\to$ GCF\_050627025.1.
  \item \textbf{Download.} \texttt{datasets download genome accession <acc> --include
        genome,protein,gff3} inside conda env \texttt{bvbrc28} on uicgpu (with
        \texttt{source \textasciitilde/env.sh} for the NCBI HTTP proxy). Packages validated.
  \item \textbf{Genome statistics} (\texttt{genome\_stats.py}): sequence length + GC\% from
        \texttt{*\_genomic.fna}; CDS/gene/tRNA feature counts from RefSeq \texttt{genomic.gff};
        protein count from \texttt{protein.faa}.
  \item \textbf{Trehalose pathway scan} (\texttt{treh2.py}): product-name regex over RefSeq/PGAP
        CDS \texttt{product=} fields (URL-decoded), classifying each of OtsA, OtsB, TreY, TreZ,
        TreS, trehalase; rollup into pathway calls (OtsA/OtsB, TreY/TreZ, TreS complete iff
        both member genes present).
  \item \textbf{ANI:} \texttt{fastANI} all-vs-all across the 3 PAMC strains plus the
        \emph{V.\ paradoxus} type strain.
  \item \textbf{Proteome comparison} (BV-BRC Proteome Comparison analogue): \texttt{makeblastdb}
        + \texttt{blastp} best-hit orthology ($\geq$30\%~id, $\geq$70\% query coverage,
        $e \leq 10^{-5}$) across the three PAMC proteomes.
  \item \textbf{LLM-judge.} Compact evidence bundle $\to$ Argo \texttt{argo:gpt-5.2} (free) for an
        independent verdict (no regex scoring).
\end{enumerate}

\section{Results vs paper}

\subsection{Genome statistics (C1--C4) --- Table 1}

\begin{longtable}{@{}llrrl@{}}
\toprule
Strain & Metric & Paper & This replication & Match \\
\midrule
PAMC28711 & Size (Mb)   & 4.32   & 4.32 (4{,}316{,}152 bp, 1 contig) & \checkmark\ exact \\
          & GC\%        & 66.00  & 65.97                              & \checkmark\ ($\Delta$0.03) \\
          & tRNA        & 46     & 46                                 & \checkmark\ exact \\
          & Genes/CDS   & 4232/4071 & 4141/4196 (4074 proteins)       & \checkmark\ (annot.\ version) \\
PAMC26660 & Size (Mb)   & 7.39   & 7.39 (7{,}388{,}698 bp, 1 contig) & \checkmark\ exact \\
          & GC\%        & 66.00  & 66.00                              & \checkmark\ exact \\
          & tRNA        & 52     & 52                                 & \checkmark\ exact \\
          & Genes/CDS   & 6919/6801 & 6901/6890 (6834 proteins)       & \checkmark\ \\
PAMC28562 & Size (Mb)   & 4.69   & 4.69 (4{,}693{,}528 bp, 1 contig) & \checkmark\ exact \\
          & GC\%        & 63.70  & 63.73                              & \checkmark\ ($\Delta$0.03) \\
          & tRNA        & 47     & 47                                 & \checkmark\ exact \\
          & Genes/CDS   & 4402/4298 & 4378/4361 (4319 proteins)       & \checkmark\ \\
\bottomrule
\end{longtable}

Genome sizes, GC\%, and tRNA counts match essentially exactly. Gene/CDS counts differ by
$<$2\% --- expected, since the paper counted from its original PGAP submission whereas we
count from the current RefSeq re-annotation of the same accession.

\subsection{Lowest GC among 19 (C2)}
Paper Table 1 lists GC\% for all 19 strains (range 63.70--69.06). Measured PAMC28562
GC = 63.73\% is the lowest of all 19. \textbf{Claim confirmed.}

\subsection{Trehalose biosynthetic pathways (C5 --- headline biological claim)}

\begin{longtable}{@{}llllcl@{}}
\toprule
Strain & OtsA/OtsB & TreY/TreZ & TreS & \# complete & Paper says \\
\midrule
PAMC28711 & \checkmark & \checkmark & \checkmark & 3 & 3 \checkmark \\
PAMC28562 & \checkmark & \checkmark & \checkmark & 3 & 3 \checkmark \\
PAMC26660 & \checkmark & --- & --- & 1 (OtsA/OtsB only) & 1 (OtsA/OtsB) \checkmark \\
\bottomrule
\end{longtable}

This is a clean, direct reproduction of the paper's central biological claim on the actual
genomes: the two 3-pathway strains and the single-pathway strain are correctly distinguished, and
PAMC26660's restriction to OtsA/OtsB is confirmed. Trehalase (degradation) product-name hits were
found in PAMC28711 and PAMC28562 but not PAMC26660, consistent with the paper's degradation-side
pattern at coarse resolution (family-level GH37 vs GH15 requires dbCAN2 --- see \S4.6).

\subsection{ANI / species distinctness (C6) --- Table 2}

\begin{longtable}{@{}lrrrrl@{}}
\toprule
Strain & Paper ANIb & Paper ANIm & Paper dDDH & This fastANI & $<$95\% threshold? \\
\midrule
PAMC28711 & 84.24\% & 85.61\% & 24.5\% & 82.2\% & \checkmark\ yes \\
PAMC26660 & 84.24\% & 88.01\% & 31.4\% & 85.6\% & \checkmark\ yes \\
PAMC28562 & 78.77\% & 84.96\% & 22.8\% & 81.3\% & \checkmark\ yes \\
\bottomrule
\end{longtable}

fastANI (k-mer/MinHash mapper) is a different algorithm from OAT/ANIb (BLAST) and ANIm
(MUMmer), so absolute values differ by 1--3\%. Qualitative conclusion is fully reproduced:
every PAMC strain sits well below the 95\% ANI species boundary versus \emph{V.\ paradoxus}.

\subsection{Proteome comparison}

\begin{longtable}{@{}lrr@{}}
\toprule
Query $\to$ subject & Shared orthologs / query total & \% \\
\midrule
PAMC28711 $\to$ PAMC26660 & 3254 / 4074 & 79.9\% \\
PAMC28711 $\to$ PAMC28562 & 3220 / 4074 & 79.0\% \\
PAMC28562 $\to$ PAMC26660 & 3504 / 4319 & 81.1\% \\
\bottomrule
\end{longtable}

Consistent with sub-species-threshold relatedness (ANI $\sim$82--86\%) and with the paper's
premise that the three strains are related but distinct.

\subsection{CAZyme content (C7)}
Not fully reproduced: \texttt{run\_dbcan} is not installed in the env, so per-strain
CAZyme-family totals (paper Table 3: PAMC28711 total 64, PAMC26660 84, PAMC28562 91; only
PAMC28711 has AAs; only PAMC28711 has both GH37 \& GH15 trehalases) were not recomputed at
family resolution. The product-name trehalase scan confirms the coarse degradation pattern but
cannot assign GH37 vs GH15. Closable by adding \texttt{run\_dbcan} + HMMER/dbCAN DBs.

\subsection{AZCL screening (C8)}
Wet-lab polysaccharide-degradation assay (Table 5) --- experimental, out of scope in silico.

\section{Verdict justification}

Every purely genomic, computationally-testable claim was checked against the actual public
RefSeq assemblies and reproduced:

\begin{itemize}[leftmargin=1.4em]
  \item Genome sizes, GC\%, and tRNA counts match essentially exactly (C1--C4);
        gene/CDS counts within annotation-version noise.
  \item The headline biological claim (C5) --- three trehalose pathways in PAMC28711/PAMC28562
        vs one in PAMC26660 --- is directly reproduced from annotation of the real genomes.
  \item The ``lowest GC of 19'' claim (C2) is confirmed.
  \item The species-distinctness conclusion (C6) is reproduced qualitatively (all ANI $<$95\%),
        with numeric offset attributable to fastANI vs the paper's OAT/ANIm/GGDC pipeline.
  \item Proteome comparison provides independent support for related-but-distinct genomes.
\end{itemize}

Only secondary/family-resolution analyses (exact CAZyme family table via dbCAN2; exact
ANIb/ANIm/dDDH; wet-lab AZCL) were not reproduced, none of which affect the paper's core
conclusions. Independent free LLM-judge (Argo \texttt{gpt-5.2}) returned REPLICATED,
Coverage 8/10, Agreement 9/10.

\section{Genuine critique}

Even though the paper's core conclusions replicate, the following are legitimate weaknesses
and framing concerns (not disqualifying, but worth flagging):

\begin{enumerate}[leftmargin=1.4em]
  \item \textbf{Product-name annotation is a weak evidence base for a ``pathway'' claim.}
        The trehalose pathway calls (both the paper's and this replication) rest on gene
        product names in PGAP/RefSeq annotation. Neither the paper nor this pass verifies
        that the identified \texttt{otsA/otsB/treY/treZ/treS} loci form intact, syntenic
        operons with correct start codons, transcription/translation signals, or that the
        proteins are enzymatically active. Presence-of-gene is not presence-of-pathway.
  \item \textbf{Adaptive/functional narrative is not tested.} The paper frames the
        3-pathway inventory in PAMC28711/PAMC28562 as an Antarctic cold/osmotic-stress
        adaptation. Neither the paper nor this replication provides comparative evidence
        (e.g.\ trehalose-pathway prevalence in non-Antarctic \emph{Variovorax}, or wet-lab
        trehalose accumulation under cold shock) that would distinguish adaptation from
        random lineage-level variation. Adaptation is a plausible hypothesis, not a
        demonstrated result.
  \item \textbf{PAMC26660's 7.39 Mb genome and its ``single pathway'' status are treated
        separately.} PAMC26660 is $\sim$60\% larger than its sister strains yet is the one
        lacking TreY/TreZ and TreS. The paper does not investigate whether the larger
        genome is due to plasmids, mobile elements, or genuine chromosomal expansion, and
        does not test whether trehalose gene loss (if it is loss) is compensated by other
        osmoprotectants. Genome-size heterogeneity within a nominal single genus deserves
        more explanation.
  \item \textbf{Type-strain and comparator selection is narrow for a ``species distinctness''
        claim.} Table~2 ANI/dDDH is reported against a handful of type strains
        (\emph{V.\ paradoxus}, \emph{V.\ beijingensis}, \emph{V.\ boronicumulans}). A modern
        replication would compare against all validly-named \emph{Variovorax} type genomes
        via TYGS + Type (Strain) Genome Server; the paper predates some newer type genomes
        that may sit closer to the PAMC strains. This replication used only
        \emph{V.\ paradoxus} NBRC 15149\textsuperscript{T} as a fastANI comparator, which is
        even narrower --- adequate for qualitative confirmation, not species assignment.
  \item \textbf{Numeric ANI differences (paper vs this pass) go unexplained by the paper.}
        The paper reports ANIb $\approx$ ANIm to within a couple of percent, whereas our
        fastANI values sit $\sim$1--3\% below ANIb for two strains and $\sim$2--3\% above
        for one. This is a known algorithm-choice effect (fastANI vs BLAST vs MUMmer) and
        does not overturn the $<$95\% conclusion, but users of the paper's exact numbers
        should not treat them as method-independent.
  \item \textbf{CAZyme claims in Table 3 are unverified here.} The paper's per-family
        CAZyme counts (including the distinctive claim that only PAMC28711 has both GH37
        and GH15 trehalases, and only PAMC28711 has AAs) rest entirely on dbCAN2 output
        that neither this pass nor a straight product-name scan can adjudicate. Full
        replication requires \texttt{run\_dbcan} + dbCAN's HMMER/DIAMOND/eCAMI databases.
        Until then, Table 3 remains a paper-internal claim, not an independently verified one.
  \item \textbf{Dedup risk is real.} There is a very-similar-sounding sibling paper from the
        same lab (Shrestha 2022 \emph{BMC Genomic Data}, PAMC28711 single-strain trehalose
        methods). Analysts must not merge or double-count. This report includes an explicit
        dedup note; downstream tooling should keep the two BVBRC entries distinct.
  \item \textbf{Wet-lab AZCL (Table 5) is unreplicable.} This is not a critique of the paper
        per se, but of any claim that this replication is ``complete'' --- the functional
        polysaccharide-degradation phenotype cannot be checked from sequence alone.
\end{enumerate}

\section{Coverage / Agreement}

\begin{itemize}[leftmargin=1.4em]
  \item \textbf{Coverage: 8/10} --- genome stats (C1--C4), lowest-GC (C2), trehalose pathways
        (C5, headline), species distinctness (C6), proteome comparison. Outstanding:
        dbCAN2 CAZyme-family table (C7), exact ANIb/ANIm/dDDH numeric reproduction,
        wet-lab AZCL (C8, not reproducible).
  \item \textbf{Agreement: 9/10} --- every tested claim agrees with the paper (exact on
        sizes/GC/tRNA/pathways; qualitative on ANI). Only sub-perfect element: fastANI
        values differ numerically (not directionally) from the paper's ANIb/ANIm, and
        CAZyme-family detail was not resolved. \textbf{No contradictions found.} All numbers
        come from real assemblies + standard tools; none fabricated.
\end{itemize}

\section{Limitations}

\begin{itemize}[leftmargin=1.4em]
  \item \texttt{fastANI} $\neq$ paper's OAT/ANIb/ANIm/GGDC; absolute ANI numbers differ 1--3\%
        (conclusion unchanged).
  \item CAZyme-family table (Table 3) not reproduced at dbCAN2 resolution; trehalase family
        (GH37 vs GH15) not sub-typed.
  \item Gene/CDS counts use current RefSeq re-annotation, not the paper's original PGAP
        submission (hence $<$2\% offsets).
  \item Only one of the three Table-2 type strains (\emph{V.\ paradoxus} NBRC 15149\textsuperscript{T})
        pulled for ANI; \emph{V.\ beijingensis} 502\textsuperscript{T} and
        \emph{V.\ boronicumulans} NBRC 103145\textsuperscript{T} not (would tighten C6, does
        not change $<$95\% conclusion).
  \item AZCL wet-lab screening (Table 5) experimental, not reproducible computationally.
\end{itemize}

\section{Open questions (brief)}

See companion file \texttt{open\_questions.json} for a machine-readable list. Highlights:
(1) verify OtsA/OtsB/TreS/TreY/TreZ synteny and intactness (not just product-name presence);
(2) test the adaptation narrative via cross-habitat \emph{Variovorax} pan-genome
comparison; (3) explain the $\sim$60\% genome-size gap PAMC26660 vs PAMC28711/PAMC28562;
(4) reproduce Table~3 at dbCAN2 family resolution (GH37/GH15/AAs); (5) reassess species
assignments against a modern TYGS type-genome set.

\section*{Verdict}
\begin{center}
\Large\textbf{REPLICATED}
\end{center}

\vspace{1em}
\hrule
\vspace{0.5em}
\small\texttt{WAVE\_RESULT set=BVBRC-47 paper=Shrestha2022-IJM-5067074 verdict=REPLICATED
dir=\textasciitilde/Dropbox/REPLICATE-PROJECT/BVBRC-47-Variovorax-PAMC-comparative-2022
one\_line=Genome sizes/GC\%/tRNA counts match exactly on real RefSeq assemblies and the
headline trehalose-pathway claim (3 pathways in PAMC28711/28562 vs 1 OtsA/OtsB in PAMC26660)
directly reproduced; ANI<95\% species distinctness confirmed via fastANI; free-endpoint
LLM-judge REPLICATED 8/10 coverage 9/10 agreement.}

\end{document}
